\documentclass[11pt,twoside,openany]{book}
\usepackage{textbook}
% generated by tools/compass.py
\newcommand{\compassstrip}[2]{%
\begin{tikzpicture}[x=1mm,y=1mm,font=\sffamily\scriptsize]
\draw[accent!50,line width=0.6pt] (-3,-5) -- (43.5,-5);
\node[text=accent,font=\sffamily\tiny] at (20.2,-8) {Vol.\,1};
\draw[accent!50,line width=0.6pt] (49.6,-5) -- (71.8,-5);
\node[text=accent,font=\sffamily\tiny] at (60.7,-8) {Vol.\,2};
\draw[accent!50,line width=0.6pt] (77.89999999999999,-5) -- (100.09999999999998,-5);
\node[text=accent,font=\sffamily\tiny] at (89.0,-8) {Vol.\,3};
\draw[accent!50,line width=0.6pt] (106.19999999999997,-5) -- (120.29999999999997,-5);
\node[text=accent,font=\sffamily\tiny] at (113.2,-8) {Vol.\,4};
\draw[black!40] (3,0) -- (5.1,0);
\pgfmathsetmacro\on{(#1<=1)&&(1<=#2)}
\ifdim\on pt>0.5pt\node[circle,fill=accent,text=white,inner sep=0pt,minimum size=5.4mm] at (0,0) {1};\else\node[circle,draw=black!45,text=black!60,inner sep=0pt,minimum size=5.4mm] at (0,0) {1};\fi
\draw[black!40] (11.1,0) -- (13.2,0);
\pgfmathsetmacro\on{(#1<=2)&&(2<=#2)}
\ifdim\on pt>0.5pt\node[circle,fill=accent,text=white,inner sep=0pt,minimum size=5.4mm] at (8.1,0) {2};\else\node[circle,draw=black!45,text=black!60,inner sep=0pt,minimum size=5.4mm] at (8.1,0) {2};\fi
\draw[black!40] (19.2,0) -- (21.299999999999997,0);
\pgfmathsetmacro\on{(#1<=3)&&(3<=#2)}
\ifdim\on pt>0.5pt\node[circle,fill=accent,text=white,inner sep=0pt,minimum size=5.4mm] at (16.2,0) {3};\else\node[circle,draw=black!45,text=black!60,inner sep=0pt,minimum size=5.4mm] at (16.2,0) {3};\fi
\draw[black!40] (27.299999999999997,0) -- (29.4,0);
\pgfmathsetmacro\on{(#1<=4)&&(4<=#2)}
\ifdim\on pt>0.5pt\node[circle,fill=accent,text=white,inner sep=0pt,minimum size=5.4mm] at (24.299999999999997,0) {4};\else\node[circle,draw=black!45,text=black!60,inner sep=0pt,minimum size=5.4mm] at (24.299999999999997,0) {4};\fi
\draw[black!40] (35.4,0) -- (37.5,0);
\pgfmathsetmacro\on{(#1<=5)&&(5<=#2)}
\ifdim\on pt>0.5pt\node[circle,fill=accent,text=white,inner sep=0pt,minimum size=5.4mm] at (32.4,0) {5};\else\node[circle,draw=black!45,text=black!60,inner sep=0pt,minimum size=5.4mm] at (32.4,0) {5};\fi
\draw[black!40] (43.5,0) -- (49.6,0);
\pgfmathsetmacro\on{(#1<=6)&&(6<=#2)}
\ifdim\on pt>0.5pt\node[circle,fill=accent,text=white,inner sep=0pt,minimum size=5.4mm] at (40.5,0) {6};\else\node[circle,draw=black!45,text=black!60,inner sep=0pt,minimum size=5.4mm] at (40.5,0) {6};\fi
\draw[black!40] (55.6,0) -- (57.7,0);
\pgfmathsetmacro\on{(#1<=7)&&(7<=#2)}
\ifdim\on pt>0.5pt\node[circle,fill=accent,text=white,inner sep=0pt,minimum size=5.4mm] at (52.6,0) {7};\else\node[circle,draw=black!45,text=black!60,inner sep=0pt,minimum size=5.4mm] at (52.6,0) {7};\fi
\draw[black!40] (63.7,0) -- (65.8,0);
\pgfmathsetmacro\on{(#1<=8)&&(8<=#2)}
\ifdim\on pt>0.5pt\node[circle,fill=accent,text=white,inner sep=0pt,minimum size=5.4mm] at (60.7,0) {8};\else\node[circle,draw=black!45,text=black!60,inner sep=0pt,minimum size=5.4mm] at (60.7,0) {8};\fi
\draw[black!40] (71.8,0) -- (77.89999999999999,0);
\pgfmathsetmacro\on{(#1<=9)&&(9<=#2)}
\ifdim\on pt>0.5pt\node[circle,fill=accent,text=white,inner sep=0pt,minimum size=5.4mm] at (68.8,0) {9};\else\node[circle,draw=black!45,text=black!60,inner sep=0pt,minimum size=5.4mm] at (68.8,0) {9};\fi
\draw[black!40] (83.89999999999999,0) -- (85.99999999999999,0);
\pgfmathsetmacro\on{(#1<=10)&&(10<=#2)}
\ifdim\on pt>0.5pt\node[circle,fill=accent,text=white,inner sep=0pt,minimum size=5.4mm] at (80.89999999999999,0) {10};\else\node[circle,draw=black!45,text=black!60,inner sep=0pt,minimum size=5.4mm] at (80.89999999999999,0) {10};\fi
\draw[black!40] (91.99999999999999,0) -- (94.09999999999998,0);
\pgfmathsetmacro\on{(#1<=11)&&(11<=#2)}
\ifdim\on pt>0.5pt\node[circle,fill=accent,text=white,inner sep=0pt,minimum size=5.4mm] at (88.99999999999999,0) {11};\else\node[circle,draw=black!45,text=black!60,inner sep=0pt,minimum size=5.4mm] at (88.99999999999999,0) {11};\fi
\draw[black!40] (100.09999999999998,0) -- (106.19999999999997,0);
\pgfmathsetmacro\on{(#1<=12)&&(12<=#2)}
\ifdim\on pt>0.5pt\node[circle,fill=accent,text=white,inner sep=0pt,minimum size=5.4mm] at (97.09999999999998,0) {12};\else\node[circle,draw=black!45,text=black!60,inner sep=0pt,minimum size=5.4mm] at (97.09999999999998,0) {12};\fi
\draw[black!40] (112.19999999999997,0) -- (114.29999999999997,0);
\pgfmathsetmacro\on{(#1<=13)&&(13<=#2)}
\ifdim\on pt>0.5pt\node[circle,fill=accent,text=white,inner sep=0pt,minimum size=5.4mm] at (109.19999999999997,0) {13};\else\node[circle,draw=black!45,text=black!60,inner sep=0pt,minimum size=5.4mm] at (109.19999999999997,0) {13};\fi
\pgfmathsetmacro\on{(#1<=14)&&(14<=#2)}
\ifdim\on pt>0.5pt\node[circle,fill=accent,text=white,inner sep=0pt,minimum size=5.4mm] at (117.29999999999997,0) {14};\else\node[circle,draw=black!45,text=black!60,inner sep=0pt,minimum size=5.4mm] at (117.29999999999997,0) {14};\fi
\end{tikzpicture}}


\newcommand{\bookpart}[6]{% numeral, title, first q, last q, volume line, intro
  \cleardoublepage\thispagestyle{empty}\refstepcounter{part}%
  \addcontentsline{toc}{part}{Part #1\quad #2}%
  \vspace*{28mm}
  {\centering\sffamily
   {\color{accent}\large #5\par}\vspace{6mm}
   {\color{accent}\fontsize{14}{16}\selectfont PART #1\par}\vspace{4mm}
   {\fontsize{26}{31}\selectfont\bfseries #2\par}\vspace{12mm}
   \compassstrip{#3}{#4}\par\vspace{3mm}
   {\small\color{quiet} The questions this part answers are shown in colour.\par}}
  \vspace{12mm}
  \noindent\begin{minipage}{\textwidth}\small #6\end{minipage}
  \clearpage}

\begin{document}
\frontmatter
\begin{titlepage}
  \centering\sffamily
  \vspace*{30mm}
  {\fontsize{30}{36}\selectfont\bfseries A Textbook of Physics\par}
  \vspace{5mm}
  {\Large From Quaternions to the Action Principle\par}
  \vspace{18mm}
  {\color{accent}\rule{60mm}{1pt}\par}
  \vspace{8mm}
  {\large Volume 1\par}\vspace{2mm}
  {\LARGE\bfseries Numbers and Equations\par}
  \vspace{30mm}
  {\large Ryojiro Kurokawa\par}
  \vfill
  {\small\color{quiet} Sample for layout review: Prologue and Part I, Chapters 1--2\\
   Draft of 13 November 2026\par}
\end{titlepage}

\tableofcontents

\mainmatter
\setcounter{chapter}{-1}
\chapter[Prologue: Three Formulas Side by Side]{Prologue:\\ Three Formulas Side by Side}
\label{ch:prologue}

\begin{question}
Some sums of squares can be written as the product of two simpler factors. Which ones? And what happens when such a product is zero although neither factor is zero? This book follows that single question from Hamilton's quaternions to the dynamics of spacetime.
\end{question}

\section{A pattern from complex numbers}

Every reader of this book has met the identity
\begin{equation}\label{eq:pro-complex}
  x^2+y^2=(x+iy)(x-iy)=z\bar z ,\qquad z=x+iy .
\end{equation}
It looks like a small piece of algebra, but two large consequences follow from it.

The first is that lengths can be multiplied. Since $\overline{zw}=\bar z\,\bar w$ and complex numbers commute,
\[
  |zw|^2=(zw)(\overline{zw})=(z\bar z)(w\bar w)=|z|^2|w|^2 .
\]
So $|zw|=|z|\,|w|$: multiplying by a complex number of length one does not change lengths. It is a rotation of the plane. Division is possible too, because $z^{-1}=\bar z/|z|^2$ whenever $z\neq0$.

The second consequence concerns derivatives. The same factorisation applied to $\partial_x$ and $\partial_y$ gives
\begin{equation}\label{eq:pro-laplace2}
  \partial_x^2+\partial_y^2=(\partial_x-i\partial_y)(\partial_x+i\partial_y).
\end{equation}
The Cauchy--Riemann equations of a complex function $f=u+iv$ say exactly that $(\partial_x+i\partial_y)f=0$. Applying the other factor shows at once that $u$ and $v$ satisfy Laplace's equation. A second-order equation has been reduced to a first-order one, and the first-order equation carries the whole theory of complex functions.

The question of this book is what happens when we try the same trick with more than two squares.

\section{Three formulas}

Here are three factorisations that the book will establish in turn. The first is Hamilton's: for a quaternion $q=a+bi+cj+dk$ and its conjugate $\bar q=a-bi-cj-dk$,
\begin{equation}\label{eq:pro-qqbar}
  q\bar q=a^2+b^2+c^2+d^2 .
\end{equation}
The second is Fueter's: the operator $D=\partial_t+i\partial_x+j\partial_y+k\partial_z$ and its conjugate $\bar D$ factor the four-dimensional Laplacian,
\begin{equation}\label{eq:pro-dbard}
  \bar D D=\partial_t^2+\partial_x^2+\partial_y^2+\partial_z^2 .
\end{equation}
The third gives the length of special relativity. Let $\iota$ be a second imaginary unit with $\iota^2=-1$ which commutes with $i$, $j$ and $k$ (we meet it properly in Chapter~6), and put $X=t+\iota(xi+yj+zk)$. Then
\begin{equation}\label{eq:pro-mink}
  X\bar X=t^2-x^2-y^2-z^2 .
\end{equation}
The same number is the determinant of the $2\times2$ matrix $t+x\sigma_1+y\sigma_2+z\sigma_3$ built from the Pauli matrices, which is the form in which physicists usually meet it.

Each of the three identities is a product of two sums of four terms, so each can be computed by filling in a $4\times4$ table. \Cref{tab:pro-qqbar,tab:pro-dbard,tab:pro-mink} show the three tables. Every cell was computed from Hamilton's rule $i^2=j^2=k^2=ijk=-1$ (Chapter~1); the box below explains how the tables are read.

\begin{tableread}
\begin{enumerate}[leftmargin=1.4em,itemsep=0.2ex,topsep=0.2ex]
\item Rows hold the \emph{left} factor, columns the \emph{right} factor. With quaternions the order matters ($ij=k$ but $ji=-k$), and this convention never changes in the book.
\item Each cell shows the product of the two units, then the product of the two entries.
\item The colour and the border of a cell show the unit of the result: \mtswatch{uonestyle}\,$1$, \mtswatch{uistyle}\,$i$, \mtswatch{ujstyle}\,$j$, \mtswatch{ukstyle}\,$k$.
\item To read off the product, add the cells of each colour. The lines under the table do exactly this.
\item Two cells carrying the same small letter cancel each other.
\end{enumerate}
\end{tableread}

\multable{qqbar}{The product $q\bar q$. The twelve off-diagonal cells cancel in pairs, and only the diagonal survives: $q\bar q=a^2+b^2+c^2+d^2$.}{tab:pro-qqbar}

\multable{dbard}{The product $\bar D D$. The table has the same pattern as \cref{tab:pro-qqbar}: mixed derivatives cancel in pairs because $\partial_x\partial_y=\partial_y\partial_x$ while $(-i)(j)=-k$ and $(-j)(i)=+k$.}{tab:pro-dbard}

\multable{minkowski}{The product $X\bar X$ for $X=t+\iota(xi+yj+zk)$. The pattern is again that of \cref{tab:pro-qqbar}, but the factor $\iota$ in the headers turns three diagonal cells negative: $X\bar X=t^2-x^2-y^2-z^2$.}{tab:pro-mink}

Put side by side, the three tables say one thing. In each of them the off-diagonal cells cancel in pairs and only the diagonal is left; that is what it means for a sum of squares to factor. The tables differ only in their headers. Replacing numbers by derivatives turns \cref{tab:pro-qqbar} into \cref{tab:pro-dbard}. Adding $\iota$ to three of the headers turns Euclidean length into the length of special relativity, \cref{tab:pro-mink}. The whole of this book can be read as the story of what these three tables, and their relatives, imply.

\section{When a product of nonzero factors vanishes}

For complex numbers, $z\bar z=0$ forces $z=0$, because a sum of squares of real numbers vanishes only when every square does. The same holds for quaternions and \eqref{eq:pro-qqbar}. It fails for \eqref{eq:pro-mink}: the number $X=1+\iota i$ has $t=x=1$ and $y=z=0$, so $X\bar X=1-1=0$ although $X\neq0$. The nonzero $X$ with $X\bar X=0$ are exactly the points of the light cone $t=|\mathbf x|$ (\cref{fig:pro-cone}). In physical terms they are the directions of light rays.

\begin{figure}[htbp]
  \centering% The light cone: X Xbar = 0 with X nonzero
\begin{tikzpicture}[x=1cm,y=1cm,>=Stealth,font=\small]
  \fill[ui!10] (-1.6,1.6) -- (0,0) -- (1.6,1.6) arc (0:-180:1.6 and 0.32) -- cycle;
  \fill[ui!10] (-1.6,-1.6) -- (0,0) -- (1.6,-1.6) arc (0:180:1.6 and 0.32) -- cycle;
  \draw[ui,thick] (-1.6,1.6) -- (1.6,-1.6);
  \draw[ui,thick] (1.6,1.6) -- (-1.6,-1.6);
  \draw[ui,thick] (0,1.6) ellipse (1.6 and 0.32);
  \draw[ui,thick] (1.6,-1.6) arc (0:-180:1.6 and 0.32);
  \draw[ui,thick,dashed] (1.6,-1.6) arc (0:180:1.6 and 0.32);
  \draw[->,black!60] (0,-2.2) -- (0,2.4) node[above] {$t$};
  \draw[->,black!60] (-2.3,0) -- (2.4,0) node[right] {$\mathbf x$};
  \fill[uj] (1.9,0.5) circle (1.6pt) node[right,font=\scriptsize] {$X\bar X<0$};
  \fill[accent] (-1.1,1.1) circle (1.6pt);
  \draw[accent,thin] (-1.1,1.1) -- (-2.0,2.05);
  \node[anchor=east,align=right,font=\scriptsize,text=accent] at (-2.0,2.1) {$X\neq0$ with $X\bar X=0$:\\a light ray, $t=|\mathbf x|$};
  \fill[uone] (0.3,0.95) circle (1.6pt);
  \draw[uone,thin] (0.3,0.95) -- (2.0,2.3);
  \node[anchor=west,font=\scriptsize,text=uone] at (2.0,2.3) {$X\bar X>0$};
\end{tikzpicture}

  \caption{The light cone. Inside it $X\bar X>0$, outside it $X\bar X<0$. On the cone a nonzero $X$ has $X\bar X=0$: a factor of a sum of squares vanishes without being zero.}
  \label{fig:pro-cone}
\end{figure}

In the language of matrices, a nonzero element of length zero is a matrix that is not invertible; its rank has dropped. This one event, a factor of a quadratic form that vanishes without being zero, comes back again and again in the book, and in each place it means something different:
\begin{itemize}[itemsep=0.2ex]
\item in Part~II, four Fueter operators with cycled coefficients share a solution, because a $4\times4$ quaternion matrix is singular;
\item in Part~III, that singularity is explained by numbers of length zero in a Clifford algebra;
\item in Part~V, the set of all singular configurations is studied as a geometric object, the central topic of the book;
\item in Part~VII, a quantum state is pure exactly when a determinant vanishes;
\item in Part~VIII, light rays and spinors are the elements of rank one.
\end{itemize}

\section{Derivatives and fields in one table}

The tables are not limited to products that cancel. \Cref{tab:pro-dq} applies $D$ to a quaternion-valued field $q=q_t+q_x i+q_y j+q_z k$, for instance the velocity of a fluid together with a fourth quantity $q_t$.

\multable{dq}{The product $Dq$. Nothing cancels, and the four colours sort the sixteen derivatives into the familiar operations of vector calculus.}{tab:pro-dq}

Write $\mathbf q=(q_x,q_y,q_z)$ for the vector part. The four lines under the table combine into
\begin{equation}\label{eq:pro-dq}
  Dq=\bigl(\partial_t q_t-\nabla\!\cdot\mathbf q\bigr)+\bigl(\partial_t\mathbf q+\nabla q_t+\nabla\times\mathbf q\bigr),
\end{equation}
where the second bracket is read as the quaternion with vector part $\partial_t\mathbf q+\nabla q_t+\nabla\times\mathbf q$. The positions in the table tell the story: the diagonal gives the divergence, the first row the time derivative, the first column the gradient, and the remaining six cells, which come in pairs with opposite signs, give the curl. If $\mathbf q$ is the velocity of a fluid, $\nabla\cdot\mathbf q$ measures compression and $\nabla\times\mathbf q$ measures rotation (Exercise~\ref{ex:pro-vortex}). In Chapter~43 the same table, with one $\iota$ added to the first row, will produce all four of Maxwell's equations at once.

\section{Fourteen questions}

The book is organised as a chain of questions. Each chapter begins with the question left open by the previous one and ends by handing on the next. \Cref{fig:pro-compass} shows the whole chain. The fourteen questions are grouped into four volumes:
\emph{Numbers and Equations} (questions 1--6), \emph{Algebra and Rank} (7--9), \emph{States and Spacetime} (10--12) and \emph{Action and Dynamics} (13--14). Question~8, the shape of the set of singular configurations, is the axis on which the book turns; it is drawn with a heavy frame. A smaller version of the figure opens every part and shows where we are.

\begin{figure}[p]
  \centering\resizebox{\textwidth}{!}{% generated by tools/compass.py
\begin{tikzpicture}[x=1mm,y=1mm,font=\sffamily]
\fill[accentlight,rounded corners=2pt] (0,-69.1) rectangle (136,7.1);
\node[anchor=north west,font=\sffamily\bfseries\footnotesize,text=accent] at (1.5,5.6) {Vol.\,1};
\node[anchor=north west,font=\sffamily\scriptsize,text=quiet,align=left] at (1.5,1.3) {Numbers and\\Equations};
\fill[accentlight,rounded corners=2pt] (0,-109.3) rectangle (136,-70.3);
\node[anchor=north west,font=\sffamily\bfseries\footnotesize,text=accent] at (1.5,-71.8) {Vol.\,2};
\node[anchor=north west,font=\sffamily\scriptsize,text=quiet,align=left] at (1.5,-76.1) {Algebra and\\Rank};
\fill[accentlight,rounded corners=2pt] (0,-149.5) rectangle (136,-110.5);
\node[anchor=north west,font=\sffamily\bfseries\footnotesize,text=accent] at (1.5,-112.0) {Vol.\,3};
\node[anchor=north west,font=\sffamily\scriptsize,text=quiet,align=left] at (1.5,-116.3) {States and\\Spacetime};
\fill[accentlight,rounded corners=2pt] (0,-177.3) rectangle (136,-150.7);
\node[anchor=north west,font=\sffamily\bfseries\footnotesize,text=accent] at (1.5,-152.2) {Vol.\,4};
\node[anchor=north west,font=\sffamily\scriptsize,text=quiet,align=left] at (1.5,-156.5) {Action and\\Dynamics};
\draw[draw=black!55,line width=0.5pt,fill=white,rounded corners=2pt] (25,-5.3) rectangle (135,5.3);
\node[anchor=west,font=\sffamily\bfseries\scriptsize] at (26.8,2.3) {1.\ Why could Hamilton build numbers in four dimensions but not in three?};
\node[anchor=west,font=\sffamily\scriptsize,text=quiet] at (26.8,-2.4) {$\to$ multiplication that keeps length: $q\bar q=|q|^2$; only dimensions 1, 2, 4, 8};
\node[anchor=east,font=\sffamily\tiny,text=accent] at (133.8,2.3) {Part I};
\draw[-{Stealth[length=1.6mm]},black!55] (80.0,-5.3) -- (80.0,-7.1);
\draw[draw=black!55,line width=0.5pt,fill=white,rounded corners=2pt] (25,-17.7) rectangle (135,-7.1);
\node[anchor=west,font=\sffamily\bfseries\scriptsize] at (26.8,-10.1) {2.\ What does a unit quaternion do?};
\node[anchor=west,font=\sffamily\scriptsize,text=quiet] at (26.8,-14.8) {$\to$ rotations of space (covered twice), the sphere of imaginary units, Hopf circles};
\node[anchor=east,font=\sffamily\tiny,text=accent] at (133.8,-10.1) {Part I};
\draw[-{Stealth[length=1.6mm]},black!55] (80.0,-17.7) -- (80.0,-19.5);
\draw[draw=black!55,line width=0.5pt,fill=white,rounded corners=2pt] (25,-30.1) rectangle (135,-19.5);
\node[anchor=west,font=\sffamily\bfseries\scriptsize] at (26.8,-22.5) {3.\ How do we compute with numbers that do not commute?};
\node[anchor=west,font=\sffamily\scriptsize,text=quiet] at (26.8,-27.2) {$\to$ left and right matrices; complexifying brings nonzero elements of length zero};
\node[anchor=east,font=\sffamily\tiny,text=accent] at (133.8,-22.5) {Part I};
\draw[-{Stealth[length=1.6mm]},black!55] (80.0,-30.1) -- (80.0,-31.9);
\draw[draw=black!55,line width=0.5pt,fill=white,rounded corners=2pt] (25,-42.5) rectangle (135,-31.9);
\node[anchor=west,font=\sffamily\bfseries\scriptsize] at (26.8,-34.9) {4.\ Can complex analysis be carried to four dimensions?};
\node[anchor=west,font=\sffamily\scriptsize,text=quiet] at (26.8,-39.6) {$\to$ the Fueter operator $D$, with $\bar D D=\Delta$};
\node[anchor=east,font=\sffamily\tiny,text=accent] at (133.8,-34.9) {Part II};
\draw[-{Stealth[length=1.6mm]},black!55] (80.0,-42.5) -- (80.0,-44.3);
\draw[draw=black!55,line width=0.5pt,fill=white,rounded corners=2pt] (25,-54.9) rectangle (135,-44.3);
\node[anchor=west,font=\sffamily\bfseries\scriptsize] at (26.8,-47.3) {5.\ What happens if the coefficients $1,i,j,k$ are cycled?};
\node[anchor=west,font=\sffamily\scriptsize,text=quiet] at (26.8,-52.0) {$\to$ a common linear solution appears: the coefficient matrix $C$ is singular (rank 3)};
\node[anchor=east,font=\sffamily\tiny,text=accent] at (133.8,-47.3) {Part II};
\draw[-{Stealth[length=1.6mm]},black!55] (80.0,-54.9) -- (80.0,-56.7);
\draw[draw=black!55,line width=0.5pt,fill=white,rounded corners=2pt] (25,-67.3) rectangle (135,-56.7);
\node[anchor=west,font=\sffamily\bfseries\scriptsize] at (26.8,-59.7) {6.\ Is that singularity an accident?};
\node[anchor=west,font=\sffamily\scriptsize,text=quiet] at (26.8,-64.4) {$\to$ no: the DFT eigenvalues are paravectors of length zero (nullity theorem)};
\node[anchor=east,font=\sffamily\tiny,text=accent] at (133.8,-59.7) {Part III};
\draw[-{Stealth[length=1.6mm]},black!55] (80.0,-67.3) -- (80.0,-72.1);
\draw[draw=black!55,line width=0.5pt,fill=white,rounded corners=2pt] (25,-82.7) rectangle (135,-72.1);
\node[anchor=west,font=\sffamily\bfseries\scriptsize] at (26.8,-75.1) {7.\ Which algebra hides in the orderings and coefficients?};
\node[anchor=west,font=\sffamily\scriptsize,text=quiet] at (26.8,-79.8) {$\to$ $S_4$ and the Klein group, Clifford systems, octonions};
\node[anchor=east,font=\sffamily\tiny,text=accent] at (133.8,-75.1) {Part IV};
\draw[-{Stealth[length=1.6mm]},black!55] (80.0,-82.7) -- (80.0,-84.5);
\draw[draw=accent,line width=1.1pt,fill=accent!12,rounded corners=2pt] (25,-95.1) rectangle (135,-84.5);
\node[anchor=west,font=\sffamily\bfseries\scriptsize] at (26.8,-87.5) {8.\ What shape does the set of singular configurations have?};
\node[anchor=west,font=\sffamily\scriptsize,text=quiet] at (26.8,-92.2) {$\to$ rank geometry: two mechanisms by handedness; odd rank 7 after complexifying};
\node[anchor=east,font=\sffamily\tiny,text=accent] at (133.8,-87.5) {Part V};
\draw[-{Stealth[length=1.6mm]},black!55] (80.0,-95.1) -- (80.0,-96.9);
\draw[draw=black!55,line width=0.5pt,fill=white,rounded corners=2pt] (25,-107.5) rectangle (135,-96.9);
\node[anchor=west,font=\sffamily\bfseries\scriptsize] at (26.8,-99.9) {9.\ Why does one formula have two real forms?};
\node[anchor=west,font=\sffamily\scriptsize,text=quiet] at (26.8,-104.6) {$\to$ $t^2+|\mathbf x|^2$ and $t^2-|\mathbf x|^2$; on the Lorentz side rank 7 is a shared light ray};
\node[anchor=east,font=\sffamily\tiny,text=accent] at (133.8,-99.9) {Part VI};
\draw[-{Stealth[length=1.6mm]},black!55] (80.0,-107.5) -- (80.0,-112.3);
\draw[draw=black!55,line width=0.5pt,fill=white,rounded corners=2pt] (25,-122.9) rectangle (135,-112.3);
\node[anchor=west,font=\sffamily\bfseries\scriptsize] at (26.8,-115.3) {10.\ What does length zero mean in quantum theory?};
\node[anchor=west,font=\sffamily\scriptsize,text=quiet] at (26.8,-120.0) {$\to$ $\det\rho=0$: pure states; entanglement of two qubits};
\node[anchor=east,font=\sffamily\tiny,text=accent] at (133.8,-115.3) {Part VII};
\draw[-{Stealth[length=1.6mm]},black!55] (80.0,-122.9) -- (80.0,-124.7);
\draw[draw=black!55,line width=0.5pt,fill=white,rounded corners=2pt] (25,-135.3) rectangle (135,-124.7);
\node[anchor=west,font=\sffamily\bfseries\scriptsize] at (26.8,-127.7) {11.\ What does length zero mean in relativity?};
\node[anchor=west,font=\sffamily\scriptsize,text=quiet] at (26.8,-132.4) {$\to$ light rays = rank one = spinors; mass is measured from the equator};
\node[anchor=east,font=\sffamily\tiny,text=accent] at (133.8,-127.7) {Part VIII};
\draw[-{Stealth[length=1.6mm]},black!55] (80.0,-135.3) -- (80.0,-137.1);
\draw[draw=black!55,line width=0.5pt,fill=white,rounded corners=2pt] (25,-147.7) rectangle (135,-137.1);
\node[anchor=west,font=\sffamily\bfseries\scriptsize] at (26.8,-140.1) {12.\ What survives when states and spacetime move?};
\node[anchor=west,font=\sffamily\scriptsize,text=quiet] at (26.8,-144.8) {$\to$ Berry phase, Thomas--Wigner rotation, curved space};
\node[anchor=east,font=\sffamily\tiny,text=accent] at (133.8,-140.1) {Part IX};
\draw[-{Stealth[length=1.6mm]},black!55] (80.0,-147.7) -- (80.0,-152.5);
\draw[draw=black!55,line width=0.5pt,fill=white,rounded corners=2pt] (25,-163.1) rectangle (135,-152.5);
\node[anchor=west,font=\sffamily\bfseries\scriptsize] at (26.8,-155.5) {13.\ Kinematics is complete. What drives it?};
\node[anchor=west,font=\sffamily\scriptsize,text=quiet] at (26.8,-160.2) {$\to$ the action principle: Maxwell, gravity, throats, collapse, quantisation};
\node[anchor=east,font=\sffamily\tiny,text=accent] at (133.8,-155.5) {Part X--XI};
\draw[-{Stealth[length=1.6mm]},black!55] (80.0,-163.1) -- (80.0,-164.9);
\draw[draw=black!55,line width=0.5pt,fill=white,rounded corners=2pt] (25,-175.5) rectangle (135,-164.9);
\node[anchor=west,font=\sffamily\bfseries\scriptsize] at (26.8,-167.9) {14.\ What do we still not know?};
\node[anchor=west,font=\sffamily\scriptsize,text=quiet] at (26.8,-172.6) {$\to$ open problems; the fine-structure constant $\alpha$ is not derived};
\node[anchor=east,font=\sffamily\tiny,text=accent] at (133.8,-167.9) {Part Epilogue};
\end{tikzpicture}
}
  \caption{The chain of fourteen questions. Each box gives a question and, after the arrow, the short form of its answer; the answer raises the next question.}
  \label{fig:pro-compass}
\end{figure}

\section{How each chapter is built}

Every chapter follows the same plan, so that a reader always knows where she is.
\begin{enumerate}[itemsep=0.2ex]
\item \textbf{The question}, in a shaded box, repeats the question left by the previous chapter.
\item The first section explains \textbf{why one would ask it}: the history, the physics that needs it, or something odd in the previous computation.
\item Next comes \textbf{a first computation}, small enough to do completely by hand.
\item Where a product sits at the centre of the chapter, it is shown in \textbf{a multiplication table}.
\item Then the general statements follow, as definitions, theorems and \textbf{proofs written to be checked by hand}. Tools from outside the main line are given in grey \textbf{toolbox} frames at the moment they are needed.
\item \textbf{Figures} are referred to in the text and are not decorative.
\item A framed box at the end says \textbf{what we learned and what comes next}.
\item \textbf{Exercises} and a short section \textbf{where this comes from} close the chapter. That section separates classical results, results of the author's research, and interpretation.
\end{enumerate}

\paragraph{What the book does not do.} The book does not derive the value of the fine-structure constant $\alpha$, and no chapter claims to. Where a result is a conjecture, or a check of a conjecture in a special case, the text says so.

\paragraph{What the reader needs.} For Volume~1: linear algebra (matrices, determinants, eigenvalues), calculus in several variables, and complex numbers. Later volumes say at their start what they add.

\begin{learned}
A sum of squares that factors gives two things: lengths that multiply, and first-order equations whose solutions are harmonic. The three formulas $q\bar q=|q|^2$, $\bar DD=\Delta$ and $X\bar X=t^2-|\mathbf x|^2$ share one multiplication table, and in the third a nonzero factor can have length zero. The first formula is the oldest. Chapter~1 asks why Hamilton needed four dimensions to obtain it.
\end{learned}

\begin{exercises}
\item Fill in the $2\times2$ multiplication table of $(\partial_x-i\partial_y)(\partial_x+i\partial_y)$ and check \eqref{eq:pro-laplace2}.
\item Check the cell in row $\iota y\,j$ and column $-\iota z\,k$ of \cref{tab:pro-mink}, and the cell with which it cancels.
\item Show that $\det(t+x\sigma_1+y\sigma_2+z\sigma_3)=t^2-x^2-y^2-z^2$, where $\sigma_1=\begin{psmallmatrix}0&1\\1&0\end{psmallmatrix}$, $\sigma_2=\begin{psmallmatrix}0&-i\\i&0\end{psmallmatrix}$, $\sigma_3=\begin{psmallmatrix}1&0\\0&-1\end{psmallmatrix}$.
\item Use \cref{tab:pro-dq} to compute $Dq$ for $q=t+xi+yj+zk$. \emph{Answer:} $Dq=-2$.
\item\label{ex:pro-vortex} A fluid turns rigidly about the $z$-axis with velocity $\mathbf q=(-y,x,0)$, and $q_t=0$. Compute $Dq$ from the table and interpret each nonzero cell. \emph{Answer:} $Dq=2k$; the only contribution is the curl, and $2$ is twice the angular velocity.
\end{exercises}

\begin{sources}
Identity \eqref{eq:pro-complex} and the Cauchy--Riemann equations are classical. Formula \eqref{eq:pro-qqbar} is due to Hamilton (1843), formula \eqref{eq:pro-dbard} to Fueter (1934--35), and the matrix form of \eqref{eq:pro-mink} goes back to Pauli's spin matrices (1927). The tables are a device of this book; every cell is computed by a short program from Hamilton's rules, so that no sign is typed by hand.
\end{sources}


\bookpart{I}{The Discovery of Quaternions and Their Foundations}{1}{3}{Volume 1 \textperiodcentered\ Numbers and Equations}{%
Part I asks three questions. Why could Hamilton build a number system in four dimensions after years of failure in three (Chapters 1 and 2)? What does a quaternion of length one do to space (Chapter 3)? And how does one compute with numbers whose products depend on the order of the factors (Chapters 4 to 6)? The answers set up everything that follows: the factorisation $q\bar q=|q|^2$, the sphere of imaginary units, and, at the end of the part, the first nonzero numbers whose length is zero.}

\chapter{In Search of Three-Dimensional Numbers}
\label{ch:three}

\begin{question}
Complex numbers turn the rotations of the plane into multiplication. Is there a number system that does the same for three-dimensional space? Concretely: can the points $a+bi+cj$ of space be multiplied so that lengths multiply, $|pq|=|p|\,|q|$?
\end{question}

\section{Why ask it}

Around 1800 complex numbers acquired a picture. Wessel (1799) and Argand (1806) drew $z=x+iy$ as an arrow in the plane, and the picture explained the rule of multiplication: writing $z=r_1e^{i\theta_1}$ and $w=r_2e^{i\theta_2}$,
\[
  zw=r_1r_2\,e^{i(\theta_1+\theta_2)} ,
\]
so lengths multiply and angles add (\cref{fig:three-complex}). Multiplying by a number of length one is a rotation. A whole branch of geometry, the rotations of the plane, became arithmetic.

\begin{figure}[htbp]
  \centering% Multiplying complex numbers multiplies lengths and adds angles
\begin{tikzpicture}[x=1cm,y=1cm,>=Stealth,font=\small]
  \draw[->,black!60] (-0.5,0) -- (4.2,0) node[below left] {$\Real$};
  \draw[->,black!60] (0,-0.4) -- (0,3.4) node[below left] {$\Imag$};
  \coordinate (O) at (0,0);
  \draw[->,very thick,uone] (O) -- ({2.2*cos(15)},{2.2*sin(15)}) node[right] {$z=r_1e^{i\theta_1}$};
  \draw[->,very thick,uj] (O) -- ({1.5*cos(45)},{1.5*sin(45)}) node[above left] {$w=r_2e^{i\theta_2}$};
  \draw[->,very thick,ui] (O) -- ({3.3*cos(60)},{3.3*sin(60)}) node[above] {$zw=r_1r_2\,e^{i(\theta_1+\theta_2)}$};
  \draw[uone] (0.7,0) arc (0:15:0.7);
  \node[uone,font=\scriptsize] at ({0.98*cos(8)},{0.98*sin(8)}) {$\theta_1$};
  \draw[uj] (1.25,0) arc (0:45:1.25);
  \node[uj,font=\scriptsize] at ({1.5*cos(31)},{1.5*sin(31)}) {$\theta_2$};
  \draw[ui] (2.1,0) arc (0:60:2.1);
  \node[ui,font=\scriptsize,anchor=west] at ({2.15*cos(42)},{2.15*sin(42)}) {$\theta_1+\theta_2$};
\end{tikzpicture}

  \caption{Multiplying complex numbers multiplies their lengths and adds their angles.}
  \label{fig:three-complex}
\end{figure}

Physics happens in three dimensions. Rigid bodies turn about axes that point anywhere in space, and the composition of two such turns is again a turn about some new axis. A number system for space, with a multiplication that multiplies lengths, would do for these rotations what complex numbers do for the plane. William Rowan Hamilton in Dublin searched for such numbers for many years. He took points of space to be \emph{triples} $a+bi+cj$, with two imaginary units $i$ and $j$, and he held on to one test, which he called the law of the moduli: the length of a product must be the product of the lengths. \Cref{fig:three-timeline} places the story in time.

\begin{figure}[htbp]
  \centering% Timeline of the discovery
\begin{tikzpicture}[x=1mm,y=1mm,font=\scriptsize]
  \draw[black!50,line width=0.8pt] (0,0) -- (119,0);
  \foreach \yr in {1750,1800,1850,1900} {\draw[black!50] ({(\yr-1742)*0.70},-0.9) -- ({(\yr-1742)*0.70},0.9);}
  \node[anchor=east,text=quiet] at (0,0) {1742};
  \node[anchor=west,text=quiet] at (119,0) {1912};
  \newcommand{\ev}[5]{\fill[#3] ({(#1-1742)*0.70},0) circle (1.1); \draw[#3] ({(#1-1742)*0.70},0) -- ({(#1-1742)*0.70},#4); \node[anchor=#5,align=center,font=\scriptsize,inner sep=1pt] at ({(#1-1742)*0.70},#4) {#2};}
  \ev{1748}{Euler 1748:\\four-square\\identity}{uone}{5}{south west}
  \ev{1799}{Wessel 1799:\\complex numbers\\as arrows}{uone}{5}{south}
  \ev{1806}{Argand 1806}{uone}{-5}{north}
  \ev{1843}{Hamilton 1843:\\$i^2=j^2=k^2=ijk=-1$}{ui}{19}{south}
  \ev{1845}{Cayley 1845:\\octonions}{uk}{-5}{north}
  \ev{1878}{Frobenius 1878:\\only $\RR,\CC,\HH$}{uj}{5}{south}
  \ev{1898}{Hurwitz 1898:\\dimensions\\1, 2, 4, 8}{uj}{19}{south}
\end{tikzpicture}

  \caption{From complex numbers to the classification of number systems.}
  \label{fig:three-timeline}
\end{figure}

\section{A first computation: why triples fail}

Here is a quick experiment with whole numbers. Both $3=1^2+1^2+1^2$ and $5=2^2+1^2+0^2$ are sums of three squares. Suppose triples could be multiplied by a formula whose coefficients are integers and which obeyed the law of the moduli. Multiplying $1+i+j$ by $2+i$ would give a triple of integers whose squared length is $3\cdot5=15$. So $15$ would be a sum of three squares of integers.

It is not. The squares not exceeding $15$ are $0,1,4,9$. If $9$ is used, the other two squares must add up to $6$, and no two of $0,1,4$ do. If $9$ is not used, three squares from $0,1,4$ add up to at most $12$. So no such formula with integer coefficients exists.

The experiment is suggestive but limited, because a multiplication rule could have non-integer coefficients. The following argument closes the gap for every rule that behaves like ordinary multiplication: it is associative and has a unit, and it contains the complex numbers.

\begin{toolbox}{two facts about determinants}
For $n\times n$ real matrices, $\det(AB)=\det A\,\det B$, and $\det(-I_n)=(-1)^n$ because $-I_n$ is diagonal with $n$ entries equal to $-1$.
\end{toolbox}

\begin{theorem}[no square root of $-1$ in odd dimension]\label{thm:three-odd}
Let $A$ be a real vector space of finite dimension $n$ with an associative bilinear multiplication and a unit $1$. If $A$ contains an element $i$ with $i^2=-1$, then $n$ is even.
\end{theorem}

\begin{proof}
Let $L_i\colon A\to A$ be left multiplication by $i$, $L_i(x)=ix$. It is linear because the multiplication is bilinear. By associativity,
\[
  L_i\bigl(L_i(x)\bigr)=i(ix)=(ii)x=-x\qquad\text{for every }x\in A ,
\]
so $L_i^2=-I_n$. Taking determinants, $(\det L_i)^2=\det(-I_n)=(-1)^n$. The left-hand side is the square of a real number, hence not negative. So $(-1)^n=+1$ and $n$ is even.
\end{proof}

\begin{corollary}\label{cor:three-triples}
There is no associative multiplication with unit on the triples $a+bi+cj$ which contains the complex numbers, that is, in which $i^2=-1$.
\end{corollary}

\begin{proof}
The triples form a real vector space of dimension $3$, which is odd.
\end{proof}

\Cref{fig:three-pairing} shows what goes wrong. Multiplication by $i$ turns the plane spanned by $1$ and $i$ through a right angle: $1\mapsto i\mapsto-1$. A map whose square is $-1$ must do this to every vector, so it has to pair up the axes. With three axes one is left over: the product $ij$ has nowhere to go. With four axes the pairing works.

\begin{figure}[htbp]
  \centering\scalebox{1.25}{% Left multiplication by i pairs the axes; an odd number leaves one axis unpaired
\begin{tikzpicture}[>=Stealth,font=\small]
  \tikzset{ax/.style={circle,draw=black!60,fill=white,minimum size=7mm,inner sep=0pt}}
  \begin{scope}
    \node[font=\sffamily\bfseries\small,text=accent] at (1.1,1.55) {three axes: $1,\ i,\ j$};
    \node[ax] (a1) at (0,0) {$1$};
    \node[ax] (a2) at (1.2,0) {$i$};
    \node[ax,draw=ui,line width=1pt] (a3) at (2.4,0) {$j$};
    \draw[->,thick,uone] (a1) to[bend left=50] node[above,font=\scriptsize] {$\times i$} (a2);
    \draw[->,thick,uone] (a2) to[bend left=50] node[below,font=\scriptsize] {$\times i$: $-1$} (a1);
    \draw[->,thick,ui,dashed] (a3) to[out=60,in=120,looseness=6] node[above,font=\scriptsize] {?} (a3);
    \node[font=\scriptsize,text=quiet,align=center] at (1.2,-1.35) {$ij$ has nowhere to go:\\$\det(L_i)^2=\det(-I_3)=-1$};
  \end{scope}
  \begin{scope}[xshift=5.4cm]
    \node[font=\sffamily\bfseries\small,text=accent] at (1.8,1.55) {four axes: $1,\ i,\ j,\ k$};
    \node[ax] (b1) at (0,0) {$1$};
    \node[ax] (b2) at (1.2,0) {$i$};
    \node[ax] (b3) at (2.4,0) {$j$};
    \node[ax] (b4) at (3.6,0) {$k$};
    \draw[->,thick,uone] (b1) to[bend left=50] node[above,font=\scriptsize] {$\times i$} (b2);
    \draw[->,thick,uone] (b2) to[bend left=50] node[below,font=\scriptsize] {$-1$} (b1);
    \draw[->,thick,uj] (b3) to[bend left=50] node[above,font=\scriptsize] {$\times i$} (b4);
    \draw[->,thick,uj] (b4) to[bend left=50] node[below,font=\scriptsize] {$-j$} (b3);
    \node[font=\scriptsize,text=quiet,align=center] at (1.8,-1.35) {two pairs, each turned by $90^\circ$:\\$\det(L_i)=1$, $\ L_i^2=-I_4$};
  \end{scope}
\end{tikzpicture}
}
  \caption{Multiplication by $i$ must turn every vector through a right angle, so it pairs the axes. Three axes cannot be paired; four can.}
  \label{fig:three-pairing}
\end{figure}

\section{The leap to four: Hamilton's rule}

On 16 October 1843, walking with his wife along the Royal Canal in Dublin, Hamilton saw the way out: the product $ij$ need not be a triple at all. It could be a new, fourth unit $k$. He cut the rule into the stone of Broom Bridge, which he passed on the walk:
\begin{equation}\label{eq:three-rule}
  i^2=j^2=k^2=ijk=-1 .
\end{equation}
Everything else follows from this line together with associativity.

\begin{proposition}\label{prop:three-six}
If the multiplication is associative and real numbers commute with $i,j,k$, then \eqref{eq:three-rule} implies
\[
  ij=k,\quad jk=i,\quad ki=j,\qquad ji=-k,\quad kj=-i,\quad ik=-j .
\]
\end{proposition}

\begin{proof}
Multiply $ijk=-1$ on the right by $k$. Since $k^2=-1$, the left side becomes $ij(k^2)=-ij$, and the right side becomes $-k$; so $ij=k$.
Multiply $ijk=-1$ on the left by $i$: $(i^2)jk=-i$, so $-jk=-i$ and $jk=i$.
Multiply $ij=k$ on the left by $i$: $i^2j=ik$, so $ik=-j$.
Multiply $jk=i$ on the left by $j$: $j^2k=ji$, so $ji=-k$.
Finally $ki=(ij)i=i(ji)=i(-k)=-ik=j$ and $kj=(ij)j=i(j^2)=-i$.
\end{proof}

The six products are easy to remember with \cref{fig:three-cycle}: going around the circle in the direction of the arrows gives a plus sign, going against them a minus sign. \Cref{tab:three-units} collects all sixteen products of units. It is the prototype of every multiplication table in this book.

\begin{figure}[htbp]
  \centering% The cyclic rule for i, j, k
\begin{tikzpicture}[>=Stealth,font=\large]
  \def\R{1.35}
  \node[circle,draw=ui,fill=ui!13,line width=0.9pt,minimum size=9mm] (i) at (90:\R) {$i$};
  \node[circle,draw=uj,fill=uj!13,line width=0.9pt,minimum size=9mm] (j) at (-30:\R) {$j$};
  \node[circle,draw=uk,fill=uk!13,line width=0.9pt,minimum size=9mm] (k) at (210:\R) {$k$};
  \draw[->,very thick,accent] (i) to[bend left=28] (j);
  \draw[->,very thick,accent] (j) to[bend left=28] (k);
  \draw[->,very thick,accent] (k) to[bend left=28] (i);
  \node[font=\small,align=center,text=quiet] at (0,0) {with the\\arrow: $+$};
  \node[font=\small,align=left,anchor=west] at (2.2,0.55) {$ij=k,\quad jk=i,\quad ki=j$};
  \node[font=\small,align=left,anchor=west] at (2.2,-0.25) {$ji=-k,\quad kj=-i,\quad ik=-j$};
  \node[font=\small,align=left,anchor=west,text=quiet] at (2.2,-1.0) {against the arrow: $-$};
\end{tikzpicture}

  \caption{The cyclic rule. Following the arrows gives the product with a plus sign, going against them gives a minus sign.}
  \label{fig:three-cycle}
\end{figure}

\multable{units}{Hamilton's multiplication table. Row times column; the colour of a cell is the unit of the product. The table is not symmetric: $ij=k$ but $ji=-k$.}{tab:three-units}

Two features of the table matter later. Its diagonal is $1,-1,-1,-1$, which is where the minus signs of \cref{tab:pro-qqbar,tab:pro-mink} come from. And it is antisymmetric off the diagonal among $i,j,k$, which is where the cancellations come from.

\section{Quaternions and their product}

\begin{definition}\label{def:three-H}
A \emph{quaternion} is an expression $q=q_0+q_1i+q_2j+q_3k$ with real $q_0,\dots,q_3$. The set of quaternions is denoted $\HH$. The number $q_0=\Real q$ is the \emph{real part} of $q$, and $\vect q=q_1i+q_2j+q_3k$ is its \emph{vector part}, which we identify with the vector $(q_1,q_2,q_3)\in\RR^3$. A quaternion with $q_0=0$ is called \emph{pure}. Quaternions are added componentwise and multiplied by expanding, using \cref{tab:three-units} for the units and letting real numbers commute with everything.
\end{definition}

The product of $p=p_0+p_1i+p_2j+p_3k$ and $q=q_0+q_1i+q_2j+q_3k$ has sixteen terms, one for each cell of \cref{tab:three-pq}.

\multable{pq}{The product of two general quaternions. The lines under the table are the four components of $pq$.}{tab:three-pq}

\begin{theorem}[scalar--vector form of the product]\label{thm:three-sv}
For $p=p_0+\vect p$ and $q=q_0+\vect q$,
\begin{equation}\label{eq:three-sv}
  pq=\bigl(p_0q_0-\vect p\cdot\vect q\bigr)+\bigl(p_0\vect q+q_0\vect p+\vect p\times\vect q\bigr).
\end{equation}
\end{theorem}

\begin{proof}
Read the four lines under \cref{tab:three-pq}. The real line is $p_0q_0-p_1q_1-p_2q_2-p_3q_3=p_0q_0-\vect p\cdot\vect q$. The $i$ line is $p_0q_1+p_1q_0+(p_2q_3-p_3q_2)$, which is the first component of $p_0\vect q+q_0\vect p$ plus the first component of $\vect p\times\vect q$. The $j$ and $k$ lines are the second and third components in the same way.
\end{proof}

So the quaternion product contains both products of vector algebra: the dot product in its real part and the cross product in its vector part. (Historically the order was the reverse: vector algebra was cut out of Hamilton's quaternions later in the nineteenth century.) For pure quaternions the formula is short enough to remember.

\begin{corollary}\label{cor:three-pure}
For pure quaternions $u$ and $v$,
\begin{equation}\label{eq:three-uv}
  uv=-\,u\cdot v+u\times v .
\end{equation}
In particular $u^2=-|u|^2$, and $uv=-vu$ exactly when $u\perp v$.
\end{corollary}

\begin{proof}
Put $p_0=q_0=0$ in \eqref{eq:three-sv}. For $v=u$ the cross product vanishes. Exchanging $u$ and $v$ keeps $u\cdot v$ and reverses $u\times v$, so $uv=-vu$ holds exactly when $u\cdot v=0$.
\end{proof}

\multable{uv}{The product of two pure quaternions: the diagonal gives $-u\cdot v$, the six off-diagonal cells give $u\times v$.}{tab:three-uv}

\begin{figure}[htbp]
  \centering% Product of two pure quaternions (oblique view)
\begin{tikzpicture}[x=1cm,y=1cm,>=Stealth,font=\small]
  \coordinate (O) at (0,0);
  \coordinate (u) at (3.0,-0.5);
  \coordinate (v) at (1.3,0.9);
  \fill[uone!10] (O) -- (u) -- ($(u)+(v)$) -- (v) -- cycle;
  \draw[black!35] (u) -- ($(u)+(v)$) -- (v);
  \draw[->,very thick,uone] (O) -- (u) node[below right] {$u$};
  \draw[->,very thick,uj] (O) -- (v) node[above left] {$v$};
  \draw[->,very thick,ui] (O) -- (0,2.4) node[above] {$u\times v$};
  \draw[black!50] (0,0.28) -- (0.25,0.32) -- (0.25,0.04);
  \node[font=\scriptsize,text=quiet] at (2.2,0.15) {area $=|u\times v|$};
  \draw[uone!70,thin] (0.75,-0.125) arc (-9.5:34.7:0.76);
  \node[font=\scriptsize,uone!80!black] at (1.05,0.25) {$\phi$};
  \node[font=\scriptsize,align=left,anchor=west] at (3.6,1.6) {$uv=-u\cdot v+u\times v$};
  \node[font=\scriptsize,align=left,anchor=west,text=quiet] at (3.6,1.05) {$u\cdot v=|u||v|\cos\phi$};
  \node[font=\scriptsize,align=left,anchor=west,text=quiet] at (3.6,0.6) {$|u\times v|=|u||v|\sin\phi$};
\end{tikzpicture}

  \caption{For pure quaternions, $uv=-u\cdot v+u\times v$: the real part measures how parallel $u$ and $v$ are, the vector part is perpendicular to both.}
  \label{fig:three-cross}
\end{figure}

\begin{example}\label{ex:three-numeric}
Let $p=1+i+j$ and $q=2+i$. \Cref{tab:three-numeric} multiplies them term by term, and adding the cells of each colour gives
\[
  (1+i+j)(2+i)=1+3i+2j-k .
\]
Theorem~\ref{thm:three-sv} gives the same answer: with $\vect p=(1,1,0)$ and $\vect q=(1,0,0)$, the real part is $1\cdot2-\vect p\cdot\vect q=2-1=1$, and the vector part is $1\cdot(1,0,0)+2\cdot(1,1,0)+(1,1,0)\times(1,0,0)=(1,0,0)+(2,2,0)+(0,0,-1)=(3,2,-1)$. We meet this product again in Chapter~2, where its length explains the failure of the experiment with $15$.
\end{example}

\multable{numeric}{A product of two particular quaternions, cell by cell.}{tab:three-numeric}

\section{Associativity}

Hamilton's rule was designed to be associative, but that has to be checked.

\begin{toolbox}{two identities of vector algebra}
For vectors $u,v,w\in\RR^3$:
$u\cdot(v\times w)=(u\times v)\cdot w$ (the triple product), and
$u\times(v\times w)=v\,(u\cdot w)-w\,(u\cdot v)$ (the ``BAC--CAB'' rule).
\end{toolbox}

\begin{theorem}\label{thm:three-assoc}
The multiplication of quaternions is associative: $(pq)r=p(qr)$ for all $p,q,r\in\HH$.
\end{theorem}

\begin{proof}
Both sides are linear in each of $p$, $q$ and $r$ separately. Writing each factor as a real part plus a vector part, it is enough to prove the identity when each factor is either a real number or a pure quaternion. If one factor is a real number $\lambda$, both sides equal $\lambda$ times the product of the other two factors, because real numbers commute with everything. So let $u$, $v$, $w$ be pure. By \eqref{eq:three-uv}, applied twice,
\begin{align*}
  (uv)w&=\bigl(-u\cdot v+u\times v\bigr)w=-(u\cdot v)\,w-(u\times v)\cdot w+(u\times v)\times w ,\\
  u(vw)&=u\bigl(-v\cdot w+v\times w\bigr)=-(v\cdot w)\,u-u\cdot(v\times w)+u\times(v\times w).
\end{align*}
The real parts $-(u\times v)\cdot w$ and $-u\cdot(v\times w)$ agree by the triple product identity. For the vector parts, the BAC--CAB rule gives
\begin{gather*}
  (u\times v)\times w=-w\times(u\times v)=-u\,(w\cdot v)+v\,(w\cdot u),\\
  u\times(v\times w)=v\,(u\cdot w)-w\,(u\cdot v).
\end{gather*}
So the vector part of $(uv)w$ is $-(u\cdot v)w-(v\cdot w)u+(u\cdot w)v$, and that of $u(vw)$ is $-(v\cdot w)u+(u\cdot w)v-(u\cdot v)w$. They are equal.
\end{proof}

\section{Order matters}

\begin{proposition}\label{prop:three-comm}
For all $p,q\in\HH$, $\;pq-qp=2\,\vect p\times\vect q$.
\end{proposition}

\begin{proof}
In \eqref{eq:three-sv}, exchanging $p$ and $q$ changes only the sign of $\vect p\times\vect q$.
\end{proof}

\begin{corollary}\label{cor:three-centre}
Two quaternions commute exactly when their vector parts are parallel (one of them may be zero). The only quaternions that commute with every quaternion are the real numbers.
\end{corollary}

\begin{proof}
The first statement is \cref{prop:three-comm}, since $\vect p\times\vect q=0$ means the vectors are parallel. If $p$ commutes with $i$ and with $j$, then $\vect p$ is parallel to $(1,0,0)$ and to $(0,1,0)$, so $\vect p=0$.
\end{proof}

The commutator of two quaternions is twice the cross product of their vector parts. This small formula will return in Chapter~3, where it becomes the rule for composing small rotations, and in Part~VII, where the same rule governs the components of spin.

\section{What Frobenius and Hurwitz settled}

Two theorems of the late nineteenth century explain why Hamilton's search had to end in four dimensions. We state them without proof; Chapter~19 returns to them.

\begin{itemize}[itemsep=0.3ex]
\item \textbf{Frobenius (1878).} Every finite-dimensional real associative algebra in which every nonzero element has an inverse is isomorphic to $\RR$, $\CC$ or $\HH$.
\item \textbf{Hurwitz (1898).} If $\RR^n$ carries a bilinear multiplication with a unit and $|xy|=|x|\,|y|$ for all $x,y$, then $n$ is $1$, $2$, $4$ or $8$. The case $n=8$ is the algebra of octonions, found by Graves in 1843 and by Cayley in 1845; it is not associative.
\end{itemize}

\Cref{thm:three-odd} is the simplest piece of these results: it already rules out every odd dimension for associative algebras that contain $\CC$.

\begin{learned}
Three-dimensional numbers that contain $\CC$ cannot be associative, because multiplication by $i$ would square to $-I_3$. Hamilton's rule $i^2=j^2=k^2=ijk=-1$ gives a four-dimensional associative algebra $\HH$ whose product contains both the dot product and the cross product. But Hamilton's real test was the law of the moduli. Does $|pq|=|p|\,|q|$ hold in $\HH$, and why does a fourth square rescue the experiment with $15$? That is the question of Chapter~2.
\end{learned}

\begin{exercises}
\item Show that a sum of three squares of integers is never $\equiv7\pmod 8$. (The squares modulo $8$ are $0$, $1$ and $4$.) Deduce again that $15$ is not a sum of three squares.
\item Using only \eqref{eq:three-rule} and associativity, compute $(i+j)^2$ and $(i+j)(i-j)$. \emph{Answer:} $-2$ and $-2k$.
\item Write $L_i$, left multiplication by $i$ on $\HH$, as a $4\times4$ matrix in the basis $1,i,j,k$. Check that $L_i^2=-I_4$ and $\det L_i=1$.
\item Compute $(1+2i)(3-j+k)$ with \cref{tab:three-pq} and with \eqref{eq:three-sv}. \emph{Answer:} $3+6i-3j-k$.
\item Show that $u^2=-1$ for every pure quaternion of length one. So $-1$ has infinitely many square roots in $\HH$; Chapter~3 is about them.
\item Show that a quaternion commutes with $i$ if and only if it has the form $a+bi$.
\item Hamilton also tried rules with $ij=0$. Take the triples with $i^2=j^2=-1$ and $ij=ji=0$. Compute $(ij)j$ and $i(jj)$, and conclude that this multiplication is not associative.
\item Use \cref{prop:three-comm} to show that $ij-ji=2k$, and check it against \cref{tab:three-units}.
\end{exercises}

\begin{sources}
Everything in this chapter is classical. Hamilton announced the quaternions in 1843 and published them from 1844 on. The geometric picture of complex numbers is due to Wessel (1799) and Argand (1806); the octonions to Graves (1843) and Cayley (1845). The classification theorems are those of Frobenius (1878) and Hurwitz (1898). The proof of \cref{thm:three-odd} by determinants is standard. The multiplication tables are generated by the program described in the Prologue.
\end{sources}

\chapter{Multiplication That Preserves Length}
\label{ch:length}

\begin{question}
Hamilton's test for a number system of space was the law of the moduli, $|pq|=|p|\,|q|$. Triples could not be made to pass it. Does it hold for quaternions? And why does a fourth square rescue the experiment with $15=3\cdot5$?
\end{question}

\section{Why ask it}

For complex numbers the law of the moduli is the source of almost everything else. It makes multiplication by a number of length one a rotation. It makes division possible, through $z^{-1}=\bar z/|z|^2$. And it rests on one identity, $z\bar z=|z|^2$, together with the rule $\overline{zw}=\bar z\,\bar w$.

For quaternions the first ingredient has already appeared in the Prologue: \cref{tab:pro-qqbar} showed $q\bar q=a^2+b^2+c^2+d^2$. The second ingredient needs care, because quaternions do not commute. This chapter checks both and draws the consequences.

\section{The conjugate}

\begin{definition}\label{def:length-conj}
The \emph{conjugate} of $q=q_0+\vect q$ is $\bar q=q_0-\vect q$. The \emph{length} (or norm) of $q$ is $|q|=\sqrt{q_0^2+q_1^2+q_2^2+q_3^2}$.
\end{definition}

Conjugation keeps the real part and reverses the vector part, so $\bar q$ is the mirror image of $q$ in the real axis (\cref{fig:length-conj}). Two formulas follow at once: $q+\bar q=2\Real q$, and $\bar{\bar q}=q$.

\begin{figure}[htbp]
  \centering% Conjugation keeps the real part and reverses the vector part
\begin{tikzpicture}[x=1cm,y=1cm,>=Stealth,font=\small]
  \draw[->,black!60] (-2.6,0) -- (2.8,0) node[right] {$\Imag\HH\cong\RR^3$};
  \draw[->,black!60] (0,-0.3) -- (0,2.6) node[above] {$\Real\HH$};
  \draw[black!30,dashed] (-2.0,1.6) -- (2.0,1.6);
  \fill[accent] (1.7,1.6) circle (1.8pt) node[above right] {$q=a+\mathbf v$};
  \fill[ui] (-1.7,1.6) circle (1.8pt) node[above left] {$\bar q=a-\mathbf v$};
  \draw[->,thick,accent] (0,1.6) -- (1.65,1.6) node[midway,below,font=\scriptsize] {$\mathbf v$};
  \draw[->,thick,ui] (0,1.6) -- (-1.65,1.6) node[midway,below,font=\scriptsize] {$-\mathbf v$};
  \fill[black] (0,1.6) circle (1.5pt) node[above left,font=\scriptsize] {$a$};
  \node[font=\scriptsize,text=quiet,align=center] at (0,-0.75) {$\bar q$ is the mirror image of $q$ in the real axis:\\$q+\bar q=2a,\qquad q\bar q=|q|^2$};
\end{tikzpicture}

  \caption{Conjugation reflects the vector part and keeps the real part.}
  \label{fig:length-conj}
\end{figure}

With complex numbers, the conjugate of a product is the product of the conjugates. With quaternions the order of the factors has to be reversed.

\begin{proposition}\label{prop:length-reverse}
For all $p,q\in\HH$, $\;\overline{pq}=\bar q\,\bar p$.
\end{proposition}

\begin{proof}
By \eqref{eq:three-sv},
$\overline{pq}=\bigl(p_0q_0-\vect p\cdot\vect q\bigr)-\bigl(p_0\vect q+q_0\vect p+\vect p\times\vect q\bigr)$.
Apply \eqref{eq:three-sv} to $\bar q=q_0+(-\vect q)$ and $\bar p=p_0+(-\vect p)$:
\[
  \bar q\,\bar p=\bigl(q_0p_0-(-\vect q)\cdot(-\vect p)\bigr)+\bigl(q_0(-\vect p)+p_0(-\vect q)+(-\vect q)\times(-\vect p)\bigr).
\]
The real parts agree. The vector part is $-q_0\vect p-p_0\vect q+\vect q\times\vect p$, and $\vect q\times\vect p=-\vect p\times\vect q$, so the vector parts agree as well.
\end{proof}

For example $\overline{ij}=\bar k=-k$, while $\bar i\,\bar j=(-i)(-j)=ij=k$ and $\bar j\,\bar i=(-j)(-i)=ji=-k$. Only the reversed order gives the right answer.

\section{A first computation in a table}

The product $q\bar q$ was computed in \cref{tab:pro-qqbar}. Here is the reason why its twelve off-diagonal cells cancel. Take the cell in row $bi$ and column $-cj$: it is $(bi)(-cj)=-bc\,ij=-bc\,k$. Its partner sits in row $cj$ and column $-bi$: it is $(cj)(-bi)=-bc\,ji=+bc\,k$. The two cells contain the same real numbers, and their units are $ij$ and $ji$, which differ only by a sign. The cells in the first row and first column cancel for a similar reason: $a(-bi)=-ab\,i$ and $(bi)a=+ab\,i$. On the diagonal, each unit meets its own negative, and $(bi)(-bi)=-b^2i^2=+b^2$. So:

\begin{proposition}\label{prop:length-qqbar}
For every quaternion $q$, $\;q\bar q=\bar q q=|q|^2$.
\end{proposition}

\begin{proof}
The equation $q\bar q=|q|^2$ is \cref{tab:pro-qqbar}, read as above. Applying it to $\bar q$ in place of $q$ gives $\bar q\,\bar{\bar q}=|\bar q|^2$, which is $\bar q q=|q|^2$.
\end{proof}

\section{The law of the moduli}

\begin{theorem}[Hamilton's law of the moduli]\label{thm:length-moduli}
For all $p,q\in\HH$, $\;|pq|=|p|\,|q|$.
\end{theorem}

\begin{proof}
By \cref{prop:length-qqbar,prop:length-reverse} and associativity (\cref{thm:three-assoc}),
\[
  |pq|^2=(pq)\,\overline{(pq)}=(pq)(\bar q\,\bar p)=p\,(q\bar q)\,\bar p=p\,|q|^2\,\bar p .
\]
The real number $|q|^2$ commutes with $p$, so the last expression equals $|q|^2\,p\bar p=|q|^2|p|^2$. Take square roots.
\end{proof}

All three ingredients are used: the conjugate of a product reverses the order, the product $q\bar q$ is a real number, and the multiplication is associative. Triples lacked the last one (\cref{cor:three-triples}).

Writing the law of the moduli out in components gives an identity between polynomials. The components of $pq$ are the four lines under \cref{tab:three-pq}.

\begin{corollary}[Euler's four-square identity]\label{cor:length-euler}
For all real numbers $p_0,\dots,p_3,q_0,\dots,q_3$,
\begin{align*}
  &(p_0^2+p_1^2+p_2^2+p_3^2)(q_0^2+q_1^2+q_2^2+q_3^2)\\
  &\quad=(p_0q_0-p_1q_1-p_2q_2-p_3q_3)^2+(p_0q_1+p_1q_0+p_2q_3-p_3q_2)^2\\
  &\qquad+(p_0q_2-p_1q_3+p_2q_0+p_3q_1)^2+(p_0q_3+p_1q_2-p_2q_1+p_3q_0)^2 .
\end{align*}
\end{corollary}

\begin{proof}
The left side is $|p|^2|q|^2$, the right side is $|pq|^2$ written in the components of \cref{tab:three-pq}.
\end{proof}

Euler had found this identity in 1748, long before quaternions, in a letter about sums of four squares. Now we can return to the experiment of Chapter~1. In \cref{ex:three-numeric} we computed
\[
  (1+i+j)(2+i)=1+3i+2j-k ,
\]
and indeed $1^2+3^2+2^2+(-1)^2=15=3\cdot5$. With three squares the product $15$ could not be written; with four it can. The fourth component, the $-k$, is exactly the product $ji$ that triples had no room for. The fourth square is not a luxury: it is where the cross product lives.

\section{Division}

\begin{corollary}\label{cor:length-inverse}
Every nonzero quaternion $q$ has an inverse,
\[
  q^{-1}=\frac{\bar q}{|q|^2},\qquad qq^{-1}=q^{-1}q=1 .
\]
Consequently $pq=0$ implies $p=0$ or $q=0$.
\end{corollary}

\begin{proof}
The two products equal $q\bar q/|q|^2$ and $\bar q q/|q|^2$, both $1$ by \cref{prop:length-qqbar}. If $pq=0$ then $|p|\,|q|=|pq|=0$ by \cref{thm:length-moduli}.
\end{proof}

Since multiplication is not commutative, there are two kinds of division. The equation $ax=b$ is solved by $x=a^{-1}b$, and $xa=b$ by $x=ba^{-1}$, and the two answers are different in general. For $a=i$ and $b=j$:
\[
  ax=b:\ x=i^{-1}j=(-i)j=-k ,\qquad xa=b:\ x=j\,i^{-1}=j(-i)=k .
\]
So the fraction $b/a$ has no meaning in $\HH$; one must say on which side one divides.

\begin{example}\label{ex:length-inverse}
For $q=1+i+j$, $|q|^2=3$ and $q^{-1}=\tfrac13(1-i-j)$. Check: $(1+i+j)(1-i-j)=1-i-j+i-i^2-ij+j-ji-j^2=1+1-k+k+1=3$.
\end{example}

\section{Quaternions of length one}

By \cref{thm:length-moduli}, the product of two quaternions of length one has length one, and so does the inverse $q^{-1}=\bar q$. The quaternions of length one therefore form a group under multiplication. As a set it is the unit sphere
\[
  S^3=\{\,q\in\HH:\ q_0^2+q_1^2+q_2^2+q_3^2=1\,\}
\]
in four-dimensional space. A three-dimensional sphere cannot be drawn directly, but it can be sliced. Fixing the real part $q_0=a$ with $-1\le a\le1$ leaves the vector parts with $|\vect q|^2=1-a^2$: an ordinary two-dimensional sphere of radius $\sqrt{1-a^2}$, which shrinks to a point at $a=\pm1$ (\cref{fig:length-s3}).

\begin{figure}[htbp]
  \centering% The unit sphere S^3 sliced by the real part
\begin{tikzpicture}[x=1cm,y=1cm,font=\small]
  \fill[uone] (0,0) circle (2pt);
  \fill[uone] (10,0) circle (2pt);
  \node[below,font=\scriptsize] at (0,-1.0) {$a=-1$};
  \node[below,font=\scriptsize] at (10,-1.0) {$a=1$};
  \node[font=\scriptsize,text=quiet] at (0,0.42) {$q=-1$};
  \node[font=\scriptsize,text=quiet] at (10,0.42) {$q=1$};
  \foreach \a/\x/\rr in {-0.8/1.35/0.51, -0.4/3.05/0.779, 0/5.0/0.85, 0.4/6.95/0.779, 0.8/8.65/0.51} {
      \shade[ball color=uone!35] (\x,0) circle (\rr);
      \draw[uone!80!black,line width=0.4pt] (\x,0) ellipse ({\rr} and {0.28*\rr});
      \node[below,font=\scriptsize] at (\x,-1.0) {$a=\a$};
  }
  \node[font=\scriptsize,text=quiet,align=center] at (5.0,1.3) {each slice $\{\,a+\mathbf v:\ |\mathbf v|^2=1-a^2\,\}$ is a 2-sphere of radius $\sqrt{1-a^2}$};
\end{tikzpicture}

  \caption{The unit quaternions $S^3$, sliced by the real part. Each slice is a sphere of radius $\sqrt{1-a^2}$; the end slices are the points $q=\pm1$.}
  \label{fig:length-s3}
\end{figure}

The middle slice, $a=0$, consists of the pure quaternions of length one. By \cref{cor:three-pure} each of them satisfies $u^2=-|u|^2=-1$. So the equation $x^2=-1$, which has two solutions in $\CC$, has a whole sphere of solutions in $\HH$ (\cref{fig:length-imag}).

\begin{proposition}\label{prop:length-roots}
The solutions of $x^2=-1$ in $\HH$ are exactly the pure quaternions of length one.
\end{proposition}

\begin{proof}
Pure quaternions of length one are solutions, as just shown. Conversely, let $x=x_0+\vect x$ with $x^2=-1$. By \eqref{eq:three-sv}, $x^2=x_0^2-|\vect x|^2+2x_0\vect x$. The vector part gives $2x_0\vect x=0$. If $\vect x=0$ then $x^2=x_0^2\ge0$, which is not $-1$; so $\vect x\neq0$ and $x_0=0$. The real part then gives $-|\vect x|^2=-1$.
\end{proof}

\begin{figure}[htbp]
  \centering% The sphere of imaginary units inside Im H (oblique view)
\begin{tikzpicture}[x=1cm,y=1cm,>=Stealth,font=\small]
  \shade[ball color=uone!12] (0,0) circle (2);
  \draw[uone!70!black,line width=0.6pt] (0,0) circle (2);
  \draw[black!40] (2,0) arc (0:-180:2 and 0.55);
  \draw[black!25,dashed] (2,0) arc (0:180:2 and 0.55);
  \draw[black!40] (0,2) arc (90:-90:0.7 and 2);
  \draw[black!25,dashed] (0,2) arc (90:270:0.7 and 2);
  \draw[->,black!55] (0,0) -- (-1.55,-1.35) node[below left] {$i$};
  \draw[->,black!55] (0,0) -- (2.7,0) node[right] {$j$};
  \draw[->,black!55] (0,0) -- (0,2.7) node[above] {$k$};
  \fill[black!70] (-1.13,-0.98) circle (1.4pt);
  \fill[black!70] (2,0) circle (1.4pt);
  \fill[black!70] (0,2) circle (1.4pt);
  \coordinate (I) at (1.15,1.25);
  \draw[ui,thick] ($-1*(I)$) -- (I);
  \fill[ui] (I) circle (2pt) node[above right] {$I$};
  \fill[ui] ($-1*(I)$) circle (2pt) node[below left] {$-I$};
  \node[font=\scriptsize,text=quiet,align=left,anchor=west] at (2.4,-1.2) {every point $I$ of this sphere\\satisfies $I^2=-1$};
\end{tikzpicture}

  \caption{The sphere of imaginary units: every pure quaternion $I$ of length one satisfies $I^2=-1$. Each such $I$, together with $1$, spans a copy of the complex plane inside $\HH$.}
  \label{fig:length-imag}
\end{figure}

Each imaginary unit $I$ gives a copy $\{x+yI\}$ of the complex numbers inside $\HH$, because $I^2=-1$ is the only rule complex arithmetic needs. These copies are the \emph{slices} of $\HH$, and they will organise much of the book: the sphere of imaginary units reappears as the space of spin states in Part~VII and as the sky of directions of light in Part~VIII.

\begin{learned}
The law of the moduli holds in $\HH$ because $q\bar q=|q|^2$, because conjugation reverses products, and because the product is associative. Writing it in components gives Euler's four-square identity, and the fourth square is where the cross product lives. Every nonzero quaternion can be inverted, from either side. The unit quaternions form a group, the three-sphere $S^3$, and the square roots of $-1$ fill a two-sphere. What does a unit quaternion \emph{do} to the space around it? Chapter~3 shows that it rotates it, and that every rotation is obtained twice.
\end{learned}

\begin{exercises}
\item Check \cref{prop:length-reverse} for $p=i$ and $q=j+k$ by computing both sides.
\item Show that $\Real(pq)=\Real(qp)$ for all $p,q\in\HH$, although $pq\neq qp$ in general.
\item Use the four-square identity to write $7\cdot 15=105$ as a sum of four squares, starting from $7=2^2+1^2+1^2+1^2$ and $15=3^2+2^2+1^2+1^2$.
\item Solve $(1+i)x=j$ and $x(1+i)=j$. \emph{Answer:} $x=\tfrac12(j-k)$ and $x=\tfrac12(j+k)$.
\item Show that $|q^{-1}|=|q|^{-1}$ and $(pq)^{-1}=q^{-1}p^{-1}$.
\item Show that the map $q\mapsto\bar q$ is not a multiplicative homomorphism of $\HH$, but that $q\mapsto |q|$ is one from the nonzero quaternions to the positive reals.
\item Find all quaternions $x$ with $x^2=-4$.
\item (The slices.) Show that for a fixed imaginary unit $I$ the set $\{x+yI:\ x,y\in\RR\}$ is closed under multiplication and inversion, and that $x+yI\mapsto x+iy$ is an isomorphism onto $\CC$.
\end{exercises}

\begin{sources}
The law of the moduli for quaternions is Hamilton's (1843). The four-square identity appears in a letter of Euler to Goldbach of 1748. The description of the square roots of $-1$ and of the slices $\{x+yI\}$ is classical; slices are the starting point of the slice analysis of Chapter~15.
\end{sources}


\end{document}
